\section{Reduction from a polynomial finite map}\label{sec:reduction}
Let $N\subset\R^r$ be compact semialgebraic, the closure of its interior. Polynomials are defined on a neighbourhood of $N$. Suppose $F^{-1}(a_0)\cap N$ is finite, nonempty, and contained in the interior of $N$, and $G=b_0$ on that fibre. Let $a(t),b(t)$ be continuous semialgebraic arcs tending to $a_0,b_0$. Put
\begin{equation}\label{eq:discriminant}
 \Delta=F(\{x\in N:\det DF(x)=0\})\cup F(\partial N).
\end{equation}
Assume $a(t)\in F(N)\setminus\Delta$ for all sufficiently small positive $t$. Choose a positive continuous semialgebraic function $r(t)\to0$ with
\[
 4r(t)<\dist(a(t),\Delta).
\]
When $\Delta$ is empty the distance restriction is omitted. The set
$\Omega_t=\{x\in N:|F(x)-a(t)|\le 2r(t)\}$ contains no critical or boundary point. Define the rational matrix and its fast derivative by
\begin{equation}\label{eq:polar-derivatives}
 A_F(x)=DG(x)DF(x)^{-1},\qquad
 M(t)=\sup_{x\in\Omega_t}\|D_x A_F(x)DF(x)^{-1}\|.
\end{equation}
The last norm is the operator norm of the derivative of a matrix-valued function. Also put
\[
 c(t)=\max\{|G(x)-b(t)|:x\in N,\ F(x)=a(t)\}.
\]
All suprema are finite and semialgebraic. This follows from compactness, nonvanishing of $\det DF$ on $\Omega_t$, and quantifier elimination.

\begin{definition}[Polar-gap conditions]\label{def:gap}
We say the marked residual satisfies the polar-gap conditions at radius $r$ if
\begin{equation}\label{eq:gap}
 \sup_{\Omega_t}\|A_F\|\le K_0,\qquad
 c(t)/r(t)\longrightarrow0,\qquad
 \eps(t):=r(t)M(t)\longrightarrow0.
\end{equation}
These are inequalities on the original polynomial map and the marked arc, not assumptions about the existence of a normal form.
\end{definition}

\begin{theorem}[Intrinsic finite-sheet reduction]\label{thm:reduction}
Under the preceding hypotheses and \eqref{eq:gap}, the solutions of $F(x)=a(t)$ can be labelled by finitely many semialgebraic arcs $x_b(t)$, $1\le b\le B$. There are inverse maps $\phi_b(U,t)$ on $|U|\le r(t)$ with $F(\phi_b(U,t))=a(t)+U$, and these maps exhaust the indicated fast tube. Set
\[
 C_b(t)=G(x_b(t))-b(t),\qquad D_b(t)=A_F(x_b(t)).
\]
Then
\begin{equation}\label{eq:derived-form}
 R(\phi_b(U,t),t)=(U,C_b(t)+D_b(t)U+E_b(U,t)),
 \quad |E_b(U,t)|\le\tfrac12 M(t)|U|^2.
\end{equation}
For $H$ in a compact subset of $\Sym_{r+q}^{++}$, let
\begin{equation}\label{eq:sheet-cost}
 q_b(H,t)=\min_{U\in\R^r}\|(U,C_b(t)+D_b(t)U)\|_H.
\end{equation}
There is a uniform $K$ such that, for all sufficiently small $t$,
\begin{equation}\label{eq:relative}
 (1-K\eps(t))\min_b q_b(H,t)
 \le \min_{x\in N}\|R(x,t)\|_H
 \le(1+K\eps(t))\min_b q_b(H,t).
\end{equation}
The distance is exactly zero if and only if some $C_b(t)=0$.

If no $C_b$ is eventually zero, write
$C_b(t)=c_b t^{\rho_b}+o(t^{\rho_b})$, $c_b\ne0$, and set
$\rho=\max_b\rho_b$. The matrices $D_b(t)$ have limits $D_b^0$, and the complete leading residual set is
\begin{equation}\label{eq:intrinsic-sheets}
 \cE_\rho=\bigcup_{\rho_b=\rho}
 \{(u,D_b^0u+c_b):u\in\R^r\}.
\end{equation}
Here $\cE_\rho$ is defined directly as the zero-time section of the closure of
$\{(t,t^{-\rho}R(x,t)):x\in N,\ t>0\}$. In particular it does not depend on the inverse-sheet construction.
\end{theorem}
\begin{proof}
The restriction of $F$ over $B(a(t),2r(t))$ is a proper local diffeomorphism with finite nonempty fibres. Properness follows by taking inverse images inside the compact set $N$; the absence of $F(\partial N)$ keeps them in its interior. A proper local diffeomorphism is a covering: disjoint inverse neighbourhoods of the finite fibre cover a sufficiently small target neighbourhood, since otherwise a sequence of additional inverse points has a compact limit in the fibre. The ball is simply connected, so every covering component maps diffeomorphically onto it.

Apply the same argument to the map $(t,x)\mapsto(t,F(x)-a(t))$ over
$\{0<t<t_0,\ |U|<2r(t)\}$. After shortening $t_0$, the arcs are analytic and the base is homeomorphic to an interval times a ball. The covering is therefore trivial. Its components and graphs are semialgebraic; alternatively their graphs are the finitely many definable continuations of the inverse fibre. This supplies the asserted labels and charts, including the closed smaller ball.

Differentiating on a sheet gives
$D_U(G\circ\phi_b)=A_F\circ\phi_b$. Its derivative has norm at most $M(t)$ by \eqref{eq:polar-derivatives}. Taylor's integral formula on the convex $U$-ball proves \eqref{eq:derived-form}.

For a direct relative comparison set $V=C_b+D_bU$. The first coordinate implies $|U|\le|(U,V)|$, and compact metric bounds imply
\[
 |\|(U,V+E_b)\|_H-\|(U,V)\|_H|
 \le K\eps(t)\|(U,V)\|_H.
\]
The constant is independent of $C_b$, including $C_b=0$. The unrestricted minimizer in \eqref{eq:sheet-cost} is
\[
 U_b^*=-(L_b^{\mathsf T}HL_b)^{-1}L_b^{\mathsf T}H(0,C_b),
 \qquad L_b=\binom{I_r}{D_b}.
\]
Since $L_b$ contains $I_r$ and $D_b$ is uniformly bounded, $|U_b^*|\le K'|C_b|=o(r)$. Thus this minimizer is a feasible trial. The trials $U=0$ have cost $O(c(t))=o(r(t))$. Outside the fast tube the norm of the first residual coordinate exceeds $r(t)$, so norm equivalence puts every global minimizer in a chart. Apply the preceding comparison there for the lower bound and at $U_b^*$ for the upper bound. This proves \eqref{eq:relative}. A zero forces $U=0$, and then forces $C_b=0$; the converse is immediate.

Every bounded semialgebraic matrix arc has a limit, and every nonzero semialgebraic vector arc has a nonzero rational-power leading vector. Uniform bounds on $D_b$ show $q_b(H,t)$ is bounded above and below by positive constant multiples of $|C_b(t)|$: use the inverse of the bounded linear shear $(U,V)\mapsto(U,V+D_bU)$. Hence \eqref{eq:relative} selects the largest order $\rho$.

For completeness we identify the entire residual set, not only its minimum. A bounded normalized residual has $U=O(t^\rho)=o(r(t))$ and hence belongs to a chart. In that chart $t^{-\rho}E_b=o(1)$. A branch with $\rho_b<\rho$ cannot have a bounded normalized second coordinate because $D_bU=O(t^\rho)$. On a branch of order $\rho$ every limit is $(u,D_b^0u+c_b)$. Conversely choosing $U=t^\rho u$ realizes each such vector. These are the two inclusions in \eqref{eq:intrinsic-sheets}.
\end{proof}

\begin{corollary}[Coordinate invariance and uniformity]\label{cor:invariance}
Feasible source reparametrizations leave $\rho$ and $\cE_\rho$ unchanged. An invertible linear change $T$ of residual coordinates sends $\cE_\rho$ to $T\cE_\rho$ and its tensor envelope to $T\cC T^{\mathsf T}$. The same conclusion holds for a smooth observation-coordinate change at the common limiting observation with invertible derivative $T$. Time is marked; replacing $t$ by $s^m$ multiplies the order by $m$.

In a semialgebraic parameter family the estimate \eqref{eq:relative} is uniform whenever the metric bounds, $K_0$, and the two limits in \eqref{eq:gap} are uniform. Formula \eqref{eq:intrinsic-sheets} is uniform on compact strata on which the stated leading-vector and slope limits are uniform. No uniformity across a vanishing gap is implicit.
\end{corollary}
\begin{proof}
The first statements follow from the definition by the feasible residual image. For a smooth observation change, the difference of its values at two observations is $TR+o(|R|)$ uniformly as both observations approach the common limit, by the integral derivative formula. Every estimate in the proof of Theorem~\ref{thm:reduction} uses only the displayed uniform constants and limits. This proves the final assertions.
\end{proof}
The distinct affine spaces in \eqref{eq:intrinsic-sheets}, after duplicates are removed, are uniquely determined by the residual set: they have the same dimension $r$, and a real affine space cannot be covered by finitely many proper affine subspaces. This assertion concerns $\cE_\rho$, not recovery of every sheet from the smaller tensor envelope.

\begin{proposition}[Finite order tests]\label{prop:order-tests}
For fixed semialgebraic data, the nonzero quantities in \eqref{eq:gap} and \eqref{eq:polar-derivatives} admit Puiseux leading orders. Thus the two strict limits are equivalent to
\[
 \ord c>\ord r,\qquad \ord r+\ord M>0,
\]
with the convention $\ord0=+\infty$. Boundedness of $A_F$ is tested entrywise on the tube supremum. The branch residues $c_b$ and limiting slopes $D_b^0$, together with these order tests, extract the sheets. Scalar order data alone do not specify them.
\end{proposition}
\begin{proof}
Graphs of the suprema and of the finite inverse fibres are semialgebraic by quantifier elimination. One-variable semialgebraic preparation gives each nonzero leading power with nonzero coefficient. Taking the quotient or product gives the displayed equivalences. Apply the same preparation to $C_b$ and $D_b$. Appendix~\ref{app:real} explains how proper real charts recover these orders without introducing inaccessible divisor faces.
\end{proof}

\begin{corollary}[Coupled weighted initial maps]\label{cor:newton}
Give $x_i$ positive rational weight $w_i$. Suppose
\[
 F(x)=f(x)+F_{>1}(x),\qquad f(t^{w_1}x_1,\ldots,t^{w_r}x_r)=t f(x),
\]
where every monomial of $F_{>1}$ has weighted degree greater than one and $f^{-1}(0)=\{0\}$. Let $\eta$ be a regular value of $f$ with a nonempty real fibre. Suppose $a(t)=t\eta+O(t^{1+\delta})$ for some $\delta>0$, every monomial of $G$ has weighted degree at least $\nu>1$, and $b(t)=O(t^\nu)$. On a sufficiently small fixed neighbourhood of zero, Theorem~\ref{thm:reduction} applies with radius $r(t)=c_0t$. In particular
\[
 c(t)=O(t^\nu),\qquad
 \sup_{\Omega_t}\|A_F\|=O(t^{\nu-1}),\qquad
 M(t)=O(t^{\nu-2}),
\]
and the relative error is $O(t^{\nu-1})$. The fast initial equations may be coupled.
\end{corollary}
\begin{proof}
A weighted sphere is compact, and $f$ is nonzero there. Weighted homogeneity therefore implies properness of $f$: its norm is bounded below by a positive constant times the weighted radius. Consequently the regular fibre at $\eta$ is finite. A sufficiently small closed target ball about $\eta$ avoids all critical values and is covered by finitely many smooth inverse charts.

Put $x_i=t^{w_i}\xi_i$. After division by $t$, the rescaled fast map converges in $C^2$ on bounded $\xi$ sets to $f$. No inverse point over that target ball can escape to infinity in $\xi$. Indeed, on a small fixed weighted neighbourhood the higher-weight part of $F$ is at most a small multiple of the lower bound for $f$, so $|F(x)|=O(t)$ implies bounded rescaled weighted radius. The inverse function theorem and compactness now give all inverse charts and uniform first and second derivatives in $\xi$. They also give a target ball of radius a fixed multiple of $t$ avoiding critical and boundary values. Choose $c_0$ small enough to meet the factor four in the gap condition.

In a rescaled target coordinate $u=U/t$, the function $G$ on these charts and its first two $u$ derivatives are $O(t^\nu)$. One derivative with respect to $U$ divides by $t$, and two divide by $t^2$. These are exactly the estimates for $A_F$ and $M$. The centre residual has the asserted bound as well. Thus $c(t)/r(t)=O(t^{\nu-1})$ and $r(t)M(t)=O(t^{\nu-1})$, proving the claim. Uniform versions follow on compact sets of the stated regular values, coefficients, and positive lower bounds.
\end{proof}

\begin{example}[Approaching a ramified discriminant]\label{ex:collision}
Take
\[
 F(x,y)=(x^2+y^2,xy),\qquad
 a(t)=(t,t/2-t^{1+\lambda}),\quad \lambda>0\text{ rational}.
\]
The real critical-value set is the union of the rays $a_1=2a_2\ge0$ and $a_1=-2a_2\ge0$. For small positive $t$ the four inverse points are given by
\[
 x+y=\sigma\sqrt{2t-2t^{1+\lambda}},\qquad
 x-y=\tau\sqrt{2t^{1+\lambda}},\qquad \sigma,\tau\in\{-1,1\}.
\]
Thus the two inverse separations have different orders and the marked arc approaches a genuine critical-value wall. Choose $r(t)=c_0t^{1+\lambda}$ with a sufficiently small positive $c_0$. If $G$ is any vector-valued homogeneous polynomial of ordinary degree $n>2+2\lambda$ and $b(t)=O(t^{n/2})$, then the polar-gap conditions hold. In fact, uniformly on the fast tube,
\begin{align*}
 \|(DF)^{-1}\|&=O(t^{-(1+\lambda)/2}),&
 \|A_F\|&=O(t^{(n-2-\lambda)/2}),\\
 M(t)&=O(t^{(n-4-3\lambda)/2}),&
 c(t)/r(t)&=O(t^{n/2-1-\lambda}).
\end{align*}
To check the second-derivative bound, differentiate $DG\,(DF)^{-1}$ once and multiply by $(DF)^{-1}$. The terms are bounded by
$\|D^2G\|\|(DF)^{-1}\|^2$ and
$\|DG\|\|D^2F\|\|(DF)^{-1}\|^3$; the latter gives the displayed, weaker bound. Hence $r(t)M(t)=O(t^{(n-2-\lambda)/2})\to0$. Theorem~\ref{thm:reduction} applies to all four competing sheets and every exact cancellation of their vector residuals. This example is not covered by requiring the target to remain a fixed relative distance from the rescaled discriminant.
\end{example}

